% GPT-6 Astra Pro, /root/sixfolds, 21 November 2026.
\subsection{HD-S25: a uniform bound for the ordinary cyclic suppliers}

Write \(G_*=|G|\) and \(g_*=g(C)\) for a connected finite-group
cover \(C\). Let \(H\le G\), let \(\chi:H\to\mu_m\) have exact
order \(m\ge3\), and put
\[
 D=C/\ker\chi,\qquad X=C/H.
\]
For every Schoen-admissible supplier, meaning that the nonzero branch
exponents of \(D\to X\) occur in pairs \(a,-a\), the primitive
multiplicity \(h>0\) is even. There is an integral signed cycle
\(Y\) of dimension \(h/2\) on \(D^h\), with class in the direct
sum of the primitive determinant lines and with nonzero projection to
each such line. For the sum \(H_D\) of degree-one point classes on
the factors, the following uniform mass bound holds:
\begin{equation}\label{eq:S25-ordinary-bound}
 \operatorname{mass}_{H_D}(Y)
 \le m^{3h/2}(h!)^2
 \le G_*^{2g_*}(g_*!)^2.
\end{equation}
Here mass means the sum of the degrees of the positive and negative
effective parts, with their displayed multiplicities. Terms with
\(h=0\) and characters of order one or two need no such supplier.

\paragraph{Geometric input and its exceptional-divisor check.}
Use Schoen, \emph{Hodge classes on self-products of a variety with an
automorphism}, Compositio Math.65 (1988), Theorem 2.0, pp.11--12,
and Lemmas 2.3--2.6, pp.14--17. In the unramified case the
nonvanishing is also stated in the paragraph after Lemma 2.2, p.13.
His construction has a linear space
\(P_0\simeq\mathbb P^{h/2}\) in an Abel--Jacobi fibre in
\(\operatorname{Sym}^hX\). The prescribed blow-ups leave its strict
transform isomorphic to \(P_0\), and the cyclic cover splits over it.
A component descends to a cycle \(Q_0\) of degree one over \(P_0\)
in the finite quotient
\[
 q:D^h\longrightarrow W_0=D^h/G_0,
 \qquad \beta:W_0\longrightarrow\operatorname{Sym}^hX,
 \qquad |G_0|=m^{h-1}h!.
\]
Here \(G_0\) is the kernel of the sum map on \((\mathbb Z/m)^h\),
extended by the symmetric group. The normal-bundle calculation gives
nonzero primitive character projections of \([Q_0]\).

The exceptional-divisor step is valid for every nonzero paired
exponent \(a\), not just for inertia of order two. Indeed, over the
intersection of a pair \(a,-a\), the exceptional divisor is a
\(\mathbb P^1\)-bundle. The two branch sections are disjoint and
all other branch divisors are pulled back from its base, by the local
blow-up of the normal-crossing ideal \((x_1,x_2)\). For a primitive
character \(\eta\) of \(\mu_m\), its eigensheaf has zero stalk on
every nontrivially ramified stratum. On every other fibre it is the
extension by zero of a rank-one local system on \(\mathbb C^*\)
with monodromy \(\eta(a)\ne1\). Its cochain complex is
\(\mathbb C\xrightarrow{1-\eta(a)}\mathbb C\), hence acyclic;
duality also gives zero compactly supported cohomology. Proper base
change therefore kills the entire primitive eigencohomology of the
exceptional preimage. This is exactly the vanishing needed in
Schoen's Lemma 2.8, pp.18--19, to descend the nonzero class. The
proper base-change input is SGA4, Expos\'e V bis, 4.1.1, also Stacks
Project, Theorem 20.18.2, tag 09V6. Thus the nonvanishing used here
does not rely on an unproved exceptional-divisor assertion in a
later draft.

\paragraph{Integral transfer and degree.}
For a general point \(x\in X\), the incidence divisor
\(D_x=x+\operatorname{Sym}^{h-1}X\) restricts to
\(\mathcal O_{P_0}(1)\): containing \(x\) is one linear condition
on the complete linear system and hence on the chosen linear subspace.
Therefore
\[
 \int_{Q_0}(\beta^*D_x)^{h/2}=1,
 \qquad q^*\beta^*D_x\equiv mH_D.
\]
Choose an irreducible component \(Z_0\) of \(q^{-1}Q_0\) and
let \(e_0\le|G_0|\) be its generic degree over \(Q_0\). The
projection formula gives
\[
 \deg_{H_D}Z_0=e_0m^{-h/2}.
\]
Use the integral effective orbit sum, retaining repetitions,
\[
 R_0=\sum_{a\in G_0}a_*Z_0.
\]
The finite-quotient transfer identity gives
\(\operatorname{cl}(R_0)=e_0q^*[Q_0]\) in rational cohomology.
This uses the quotient transfer, not a flat pullback of a cycle through
the branch locus. A finite complex quotient of a smooth variety is an
oriented rational homology manifold, and the transfer identity
\(q^*q_*\alpha=\sum_a a_*\alpha\) follows by averaging and
Poincar\'e duality. Equivalently it can be checked on invariant
triangulated chains; stabilizer orders give precisely the repetitions
in the displayed orbit sum. Rational pullback is injective, since
\(q_*q^*=|G_0|\). Thus \(R_0\) retains every primitive
nonvanishing component, and
\[
 \deg_{H_D}R_0
 \le |G_0|^2m^{-h/2}.
\]

Let \(\delta(t)\) act on the first curve factor by the residual
cyclic action, and let
\[
 c_m(t)=\sum_{u\in(\mathbb Z/m)^\times}
              \exp(-2\pi iut/m)
\]
be the integral Ramanujan sum. Define the actual signed cycle
\[
 Y=\sum_{t=0}^{m-1}c_m(t)\delta(t)_*R_0.
\]
Its class is \(m\) times the rational projector to the primitive
residual characters. Kernel-of-sum invariance forces the character
on every curve factor to be the same; for a primitive character it
is nontrivial, so every factor contributes \(H^1(D)\), not
\(H^0(D)\) or \(H^2(D)\). Symmetric-group invariance, including
the Koszul sign in degree one, is the alternating tensor
\(\bigwedge^h H^1(D)_\eta\). This is precisely its determinant
line. Hence the cycle has the asserted pure determinant class.

The automorphisms preserve \(H_D\), and
\(|c_m(t)|\le\varphi(m)\). Consequently
\[
 \operatorname{mass}_{H_D}(Y)
 \le m\varphi(m)|G_0|^2m^{-h/2}
 \le m^2m^{2h-2}(h!)^2m^{-h/2}
 =m^{3h/2}(h!)^2.
\]
The rational primitive summand has dimension \(\varphi(m)h\), so
\(h\le g(D)\le g_*\), because \(\varphi(m)\ge2\).
Also \(m\le G_*\). These inequalities give the final bound in
\eqref{eq:S25-ordinary-bound}. The proof works with the actual cycles
on the curve powers and introduces no inverse Lefschetz operation.

For the E8 tuple of HD-S22 one may therefore use the single integer
\[
 S_{\rm ord}
 =(1393459200)^{6154444802}(3077222401!)^2
\]
for every ordinary admissible row in a verified supplier identity.
The other rows require their separate product construction.

\paragraph{A uniform bound for the actual divisor contractions.}
There is also a bound that applies uniformly to all quotient-curve
factors occurring in these suppliers. For a quotient \(p:C\to D\)
of degree \(t\), use the actual map
\[
 D\longrightarrow J(D)\xrightarrow{p^*}J(C).
\]
The norm-pullback identity \(p_*p^*=[t]\) and its polarized adjunction
give \((p^*)^*\theta_C=t\theta_D\). Equivalently, the equality
follows by integrating pulled-back one-forms on the degree-\(t\)
cover; branch points have measure zero. Its theta degree on the curve
is therefore \(t g(D)\). Riemann--Hurwitz gives
\[
 t g(D)=g_*-1+t-\tfrac12\deg R_p\le g_*-1+G_*.
\]
Let \(A\subset J(C)\) be an actual abelian subvariety of dimension
six with restricted theta polarization \(\theta_A\). On a product
of quotient curves and \(A\), each divisor pulled back from
\(2\Theta_C\) by one of the maps \(x_j\), \(x_j+g x_k\),
\(g\in G\), is globally generated and is bounded in the nef order
by
\[
 \lambda(H_{\rm curves}+\theta_A),\qquad
 \lambda=8(g_*+G_*).
\]
Indeed, the invariant positive theta form satisfies
\(\theta(x+y)\le2\theta(x)+2\theta(y)\), by the squared-norm
inequality. Multiplication by 2 for \(2\Theta_C\) gives coefficients
4 at two different factors. Repeated factors give coefficient at
most 8. Deck transformations preserve \(\Theta_C\), and the
\(A\)-inclusion pulls it back to \(\theta_A\), so the preceding
curve-degree bound proves the displayed nef inequality.

For an integral signed middle cycle \(Y\) on \(\ell\) such curve
factors, put \(k=\ell/2+3\). Intersect
\(Y\times A\) with \(k\) general members of any one of these
globally generated divisor products and push to \(A\). The resulting
signed effective decomposition has theta-degree mass at most
\begin{equation}\label{eq:S25-contraction-bound}
 \lambda^k\binom{k}{3}\,
       \operatorname{mass}_{H_{\rm curves}}(Y)\,6!\Delta_A.
\end{equation}
Properness follows by choosing the members successively outside the
finitely many loci that contain a component. The projection formula
and nef monotonicity reduce the degree to
\((H_{\rm curves}+\theta_A)^k\theta_A^3\) on \(Y\times A\).
Only its term of curve degree \(\ell/2\) and abelian degree 6
survives, giving exactly the binomial coefficient above.
For the primitive E8 factor, \(\Delta_A\le D^6\) by HD-S23.
This estimates every actual monomial contraction; any nonvanishing
invariant-transfer combination is established by the separate supplier
identity and its eigenline pairings.

\paragraph{Audit.}
This is a new full author read of the displayed proof, GPT-6 Astra Pro,
21 November 2026, inherited effort. The degree argument is a
restatement with explicit scope of the previously checked SQ4 argument;
the extension from an order-two inertia pair to an arbitrary nonzero
pair is proved above by its actual rank-one local system. The quotient
Jacobian map, norm-pullback polarization identity, nef coefficient 8,
and the exact surviving binomial term in the contraction bound were
also checked in full.
